% SOURCE_STATEMENT: "1. Why do you expect them to arise as monopole moduli spaces?\n\nAnswer: (Cherkis) string theory argument. 2. Are all 8-dimensional hyperk¨ ahler manifolds (non-compact, with some sort of well-behaved asymptotic behaviour) moduli spaces?\n\nAnswer: (Cherkis) yes, from string theory intuition. 3. What about dimension > 8? 4. One way to get gravitational instantons is by solving the vacuum Einstein equation in Lorentz space, and then performing a wick rotation (meaning t → it ). This worked in the case of Lorentzian Taub-NUT to get the Riemannian Taub-NUT. Which other instantons arrise this way?\n\nAnswer: In order to get a Riemannian metric, one needs t → t + c (time translation) be an isometry. Therefore it does work for Ak-cases but not for the others ( Dk, E6, E 7, E 8). 5. Can they arrive as \"non-trivial\" hyperk¨ ahler reductions of finite dimensional linear spaces?\n\nAnswer: Yes, for ALE spaces (Kronheimer's construction). Yes, for ALF if certain (Dancer) spaces are included with the linear spaces. 6. Is it true that every gravitational instanton has an asymptotic local triholomorphic isometry?\n\nReason: If you have a local R-action at infinity, then the metric is given in terms of a harmonic function, which arises from the hyperk¨ ahler moment map associated to this action (see the work of Gibbons). 7. When is it true that {Solutions of the Einstein equation } > {Self-Dual metrics }?\n\nReason: In general ≥.\n\nAnswers:\n\n(1) In the case of ALF the answer is >, since the Euclidean Schwarzschild metric is not self-dual. (2) Nakajima conjectured that in dimension 4, ALE and Ricci-flat implies self-dual. (3) In dimension > 4 and even, the conjecture is that ALE and Ricci-flat implies K¨ ahler. (4) (Carron) For odd dimensional manifolds an ALE Ricci-flat metric has to be flat since (a) by the Cheeger-Gromoll theorem this manifold has only one end (if not it splits isometrically as\n\nR × N with N compact); (b) The topology at infinity must be of the type ( R, ∞) × S2n/G where G is a finite subgroup of\n\nO(2 n + 1) acting freely on S2n. If g ∈ SO (2 n + 1) it must have 1 as an eigenvalue hence the only such G is\n\n{I, −I}; but the quotient is the real projective space, and there is no odd dimensional compact manifold whose boundary is the real projective space. Therefore G is trivial and the Bishop-Gromov inequality 2implies that M is the Euclidean space. Recall that the Bishop-Gromov inequality states that in a manifold\n\nM n with non-negative Ricci curvature r 7 → vol B(x, r )/w nrn is decreasing with equality everywhere if and only if M n is the Euclidean space. This ratio goes to 1 when r → 0, and in the ALE case we have that this ratio goes to 1 /card G when r → ∞.\n\nQuestions [of S. Cherkis]:"
% Preparation and final self-audit: CONTRIBUTING.md prompts C and D; source/audit artifacts retained in the local physics_statements_2026_10_11 batch.
% !TeX program = lualatex
\documentclass[11pt]{article}
\usepackage[margin=25mm]{geometry}
\usepackage{fontspec}
\setmainfont{Latin Modern Roman}
\usepackage{amsmath,amssymb,mathtools}
\usepackage{unicode-math}
\setmathfont{Latin Modern Math}
\usepackage{xurl,hyperref}
\hypersetup{colorlinks=true,urlcolor=blue}
\setcounter{secnumdepth}{0}
\setlength{\parindent}{0pt}
\setlength{\parskip}{5pt}
\setlength{\emergencystretch}{3em}
\newcommand{\sref}[2]{\hyperlink{source:#1:#2}{[S#2]}}
\newcommand{\cataloguescope}[1]{\paragraph{Recorded target.}#1}
\newcommand{\R}{\mathbb R}
\newcommand{\C}{\mathbb C}
\newcommand{\N}{\mathbb N}
\newcommand{\Z}{\mathbb Z}
\newcommand{\Q}{\mathbb Q}
\newcommand{\F}{\mathbb F}
\newcommand{\eps}{\varepsilon}
\begin{document}
% UnsolvedMath ID 20002522; source catalogue code AIM-PHYSICS-0030
\setcounter{section}{20002521}
\section{Why gravitational instantons arise from monopoles}\label{20002522}
{\small UnsolvedMath ID 20002522 · Mathematical Physics}
\subsection{Definitions and mathematical statement}

\paragraph{Shared notation.}$\mathbb N=\{1,2,\ldots\}$; the standard number fields have their usual meanings. All additional conventions are stated locally.


\paragraph{Definitions and notation.}A hyper-Kähler metric has three parallel complex structures satisfying the quaternionic relations. In four dimensions a gravitational instanton is a complete noncompact Ricci-flat geometry in the source’s self-dual setting. ALE and ALF describe different asymptotic geometries; ALF has an asymptotic circle direction over a three-dimensional base. A monopole moduli space parametrizes gauge-equivalence classes of solutions of a Bogomolny equation with specified charge and boundary behavior. A hyper-Kähler reduction takes a level set of three moment maps and quotients by the acting group. A triholomorphic isometry preserves all three complex structures. Wick rotation is a complex change of a time coordinate, requiring geometric conditions to yield a real Riemannian metric.

\paragraph{Catalogue wording retained for traceability.}
Retain the seven source questions: explain why gravitational instantons should arise as monopole moduli spaces; determine whether suitable eight-dimensional noncompact hyper-Kähler manifolds are moduli spaces and what happens above dimension eight; determine which instantons arise by Wick rotation beyond the Taub--NUT example; find nontrivial hyper-Kähler reductions of finite-dimensional linear spaces; decide whether every gravitational instanton has an asymptotic local triholomorphic isometry; and identify settings in which Ricci-flat Einstein solutions strictly exceed self-dual solutions. The source discusses ALF Euclidean Schwarzschild as a non-self-dual example, the ALE Ricci-flat/self-dual conjecture in dimension four, a Kähler conjecture in higher even dimensions, and a flatness argument in odd dimensions.
\paragraph{Scope and source qualification.}The workshop contains physical intuition and partial answers alongside questions. They are source context, not independent proofs checked here. Well-behaved asymptotics and the higher-dimensional replacement for self-duality require precise specification; the original informal breadth is preserved. The catalogue title/status was paired with a machine-generated research attempt; that attempt is not peer-reviewed evidence of settlement. The statement covers the original source questions rather than only the title’s subcase.
\subsection{Short English statement}
Understand moduli-space origins and constructions of gravitational instantons, their asymptotic symmetries, and distinctions among Ricci-flat and self-dual geometries.
\subsection{Sources}
\begin{itemize}
\item[\hypertarget{source:20002522:1}{S1}] ULAM.AI and the UnsolvedMath Contributors, UnsolvedMath, record 20002522 (AIM-PHYSICS-0030). Catalogue curation under CC BY 4.0; definitions and English exposition added locally.
\url{https://huggingface.co/datasets/ulamai/UnsolvedMath}.
\item[\hypertarget{source:20002522:2}{S2}] American Institute of Mathematics, $L^2$ harmonic forms in geometry and string theory; notes by Anda Degeratu and Mark Haskins, March 2004, 17 March 2004, Self-dual Gravitational Instantons questions.
\url{https://aimath.org/WWN/l2harmonic/l2harmonic.pdf}.
\end{itemize}
\label{end:20002522}
\end{document}
